% Source: https://artofproblemsolving.com/community/c3h1668349_triangles_with_integral_side
% Stable identity: aops-topic-1668349-q01
A triangle counted by $t_{20}$ but not by $t_{19}$ is one whose largest side is $20$. Let its other sides be $a<b<20$.

\begin{asy}
usepackage("amsmath");
size(5cm);
defaultpen(fontsize(10pt));
pair A=(0,0),B=(20,0),C=(7.6,6.2);
draw(A--B--C--cycle,linewidth(0.9));
label("$20$",(A+B)/2,S);label("$a$",(A+C)/2,NW);label("$b$",(B+C)/2,NE);
label("\mbox{}",(min(currentpicture,user=true)+max(currentpicture,user=true))/2,black); // keeps the text after the figure black when LaTeX typesets the labels inline
\end{asy}

By the triangle inequality, three lengths $a<b<20$ form a triangle exactly when $a+b>20$. For each $b$, $a$ ranges over $21-b\le a\le b-1$, which is nonempty only for $b\ge11$ and has $2b-21$ values. So
\[
t_{20}-t_{19}=\sum_{b=11}^{19}(2b-21)=1+3+5+\cdots+17=9^2=\boxed{81}.
\]
(If equal sides were allowed, the same count with $a\le b\le20$ would give $2+4+\cdots+20=110$.)
